\chapter{Cystectomy}

\section*{Introduction \& Clinical Indications}
Cystectomy is a significant surgical intervention that involves the complete or partial excision of the urinary bladder. The most common form is radical cystectomy, which entails the en bloc removal of the bladder along with adjacent organs and pelvic lymph nodes, primarily indicated for muscle-invasive bladder cancer (MIBC). This complex procedure necessitates the establishment of a new urinary diversion pathway to ensure continued renal function and maintain the patient's quality of life. Cystectomy is a cornerstone in the management of advanced bladder malignancies and select benign conditions, representing a critical intervention with profound implications for patient recovery and long-term care. The surgery requires meticulous preoperative planning, intraoperative precision, and comprehensive postoperative management to optimize outcomes.

\begin{itemize}
  \item Bladder removal surgery
  \item Radical cystoprostatectomy (in males, when the prostate and seminal vesicles are also excised)
  \item Anterior pelvic exenteration (in females, when the uterus, ovaries, fallopian tubes, and a portion of the vagina are removed)
  \item Partial cystectomy (for localized, non-invasive tumors or select benign conditions)
  \item Total cystectomy (often synonymous with radical cystectomy, particularly when referring to complete bladder removal)
  \item Cystourethrectomy (if the urethra is also excised, often for extensive urethral involvement or carcinoma in situ)
\end{itemize}

The primary surgical principle of radical cystectomy is the complete extirpation of the bladder and surrounding tissues that may harbor microscopic disease, thereby achieving oncologic control. This involves ensuring a wide surgical margin around the tumor and performing a comprehensive pelvic lymph node dissection (PLND) to remove regional lymph nodes, which are common sites for metastatic spread. In males, this typically includes the prostate and seminal vesicles, while in females, the uterus, ovaries, fallopian tubes, and a portion of the anterior vaginal wall may be excised to ensure adequate oncologic margins, particularly for tumors involving the bladder neck or trigone.

The rationale behind this extensive resection is to eliminate all gross and microscopic disease to prevent recurrence and improve long-term survival. Following bladder removal, the primary technical goal is to restore urinary continuity and function through a urinary diversion. The choice of diversion depends on patient factors, renal function, tumor characteristics, and surgeon preference, aiming to minimize complications and optimize quality of life. Common diversion types include the ileal conduit (Bricker conduit), an incontinent diversion where urine drains continuously into an external ostomy bag; and orthotopic neobladder reconstruction, where a segment of bowel is fashioned into a bladder substitute and anastomosed to the urethra, allowing for voluntary voiding. Continent cutaneous diversions, which involve an internal reservoir that is catheterized intermittently, are also an option for select patients. The meticulous execution of both the ablative and reconstructive phases is critical for successful outcomes, balancing oncologic efficacy with functional preservation.

The history of cystectomy dates back to the late 19th and early 20th centuries, characterized by high morbidity and mortality due to the lack of effective urinary diversion techniques. Early attempts at bladder removal were fraught with complications such as uremia and infection, limiting the procedure's applicability. The modern era of cystectomy began in the mid-20th century, marked by significant advancements in surgical technique and postoperative care.

In 1950, Dr. Eugene Bricker introduced the ileal conduit, a revolutionary technique that provided a safe and reliable method for urinary diversion, significantly improving the feasibility and outcomes of radical cystectomy. This innovation transformed cystectomy into a viable curative option for bladder cancer. Subsequent decades saw refinements in surgical technique, including extended pelvic lymph node dissections to improve oncologic staging and control, and the development of nerve-sparing techniques to preserve sexual function. The 1980s and 1990s marked the advent and widespread adoption of orthotopic neobladder reconstruction, pioneered by surgeons like Dr. Guido Zincke and Dr. Urs Studer, offering patients the potential for continent, urethral voiding and improved quality of life compared to external stomas. More recently, the integration of minimally invasive techniques, such as laparoscopic and robotic-assisted radical cystectomy, has further evolved the procedure, aiming to reduce surgical morbidity, blood loss, and hospital stay while maintaining oncologic efficacy.

The clinical indications for radical cystectomy include:

\begin{itemize}
  \item Muscle-invasive bladder cancer (MIBC), defined as T2, T3, or T4a disease, which is the most common indication and often requires neoadjuvant chemotherapy.
  \item High-grade non-muscle invasive bladder cancer (NMIBC) that is refractory to intravesical therapy (e.g., BCG-unresponsive carcinoma in situ or recurrent high-grade T1 disease), where progression risk is high.
  \item Recurrent, multifocal high-grade NMIBC that is not amenable to repeated transurethral resection of bladder tumor (TURBT) or intravesical treatments.
  \item Extensive benign bladder destruction due to severe interstitial cystitis, neurogenic bladder dysfunction, radiation cystitis, or chronic infection, when conservative measures have failed and quality of life is severely impacted.
  \item Bladder cancer with associated prostatic stromal invasion in men (T4a disease), which carries a high risk of local recurrence.
  \item Select cases of persistent or recurrent bladder cancer after prior definitive radiation therapy, where salvage cystectomy may be considered.
  \item Rarely, for certain non-urothelial bladder malignancies such as sarcoma, adenocarcinoma, or squamous cell carcinoma, depending on stage and resectability.
  \item Bladder trauma or severe congenital anomalies not amenable to reconstructive surgery.
\end{itemize}

Radical cystectomy is primarily considered a definitive, curative treatment for muscle-invasive bladder cancer (MIBC). In this context, it serves as the primary surgical intervention aimed at complete tumor eradication and is often combined with neoadjuvant or adjuvant chemotherapy to maximize oncologic outcomes. For high-risk non-muscle invasive bladder cancer (NMIBC) that has failed conservative management, such as intravesical Bacillus Calmette-Gu\'erin (BCG) therapy, cystectomy transitions to a secondary, or salvage, procedure. In these scenarios, it is performed to prevent progression to MIBC and to achieve disease control when less invasive options have been exhausted. For benign conditions, it is typically a last-resort secondary procedure after all other medical and less invasive surgical treatments have failed to alleviate severe symptoms and significantly impact a patient's quality of life. The decision for cystectomy is always made after careful consideration of the patient's overall health, life expectancy, and preferences, often in a multidisciplinary tumor board setting.

\section*{Relevant Surgical Anatomy}
The urinary bladder is a muscular, distensible organ located in the anterior pelvis, posterior to the pubic symphysis. Its superior surface is covered by peritoneum, while the anterior and lateral surfaces are extraperitoneal. In males, the bladder is situated superior to the prostate and anterior to the rectum, with the seminal vesicles and vas deferens located posteriorly. In females, the bladder is positioned anterior to the uterus and vagina. The ureters enter the bladder posterolaterally, and the urethra exits inferiorly. 

Key anatomical relations include the obturator nerves and vessels, external and internal iliac vessels, and hypogastric nerves, all in close proximity to the bladder and pelvic lymph nodes. 

The arterial supply to the bladder is primarily from the superior and inferior vesical arteries, branches of the internal iliac arteries. Venous drainage occurs via the vesical venous plexus, which drains into the internal iliac veins. Lymphatic drainage follows the arterial supply, primarily to the external iliac, internal iliac, and obturator lymph nodes, subsequently draining to the common iliac and para-aortic nodes. The autonomic nerve supply originates from the pelvic splanchnic nerves (parasympathetic) and hypogastric nerves (sympathetic), forming the inferior hypogastric plexus, crucial for bladder function and sexual health. 

Surgical landmarks include the obliterated umbilical arteries, the vas deferens in males, the round ligaments in females, and the ureters as they cross the common iliac vessels.

\section*{Preoperative Workup \& Preparation}
Thorough preoperative preparation is crucial for minimizing complications and optimizing outcomes in patients undergoing radical cystectomy. This typically involves:

\begin{itemize}
  \item **Comprehensive Medical Evaluation:** A detailed history and physical examination, including assessment of cardiac, pulmonary, and renal function. This often involves an electrocardiogram (ECG), chest X-ray, pulmonary function tests, and potentially a cardiology consultation for high-risk patients.
  
  \item **Laboratory Tests:** Complete blood count, comprehensive metabolic panel (including electrolytes, renal and liver function tests), coagulation profile, and urinalysis with culture to rule out infection. Blood typing and cross-matching are essential due to potential blood loss.
  
  \item **Imaging Studies:** Preoperative staging with cross-sectional imaging (CT scan of the chest, abdomen, and pelvis with intravenous contrast) to assess for metastatic disease, evaluate renal anatomy, and identify any hydronephrosis. Bone scan may be performed if clinically indicated (e.g., bone pain, elevated alkaline phosphatase).
  
  \item **Anesthesia Consultation:** Evaluation by an anesthesiologist to assess fitness for prolonged general anesthesia and discuss pain management strategies, including epidural placement.
  
  \item **Nutritional Optimization:** Assessment of nutritional status and intervention if malnutrition is present, as this can significantly impact wound healing, immune function, and overall recovery. Oral nutritional supplements may be prescribed.
  
  \item **Bowel Preparation:** Mechanical bowel preparation (e.g., polyethylene glycol solution) combined with oral antibiotics (e.g., neomycin and metronidazole) is often prescribed to reduce bacterial load in the bowel, especially for diversions involving bowel segments, to minimize infectious complications.
  
  \item **Stoma Site Marking:** For patients undergoing incontinent diversions (ileal conduit) or continent cutaneous diversions, a specialized stoma nurse will preoperatively mark an optimal stoma site on the abdomen, considering patient anatomy, body habitus, clothing lines, and ease of self-care.
  
  \item **Medication Adjustment:** Review and adjustment of medications, particularly anticoagulants or antiplatelet agents, which may need to be held or bridged prior to surgery. Diabetic medications may also require adjustment.
  
  \item **Informed Consent:** Detailed discussion with the patient and family regarding the procedure, potential risks, benefits, alternative treatments, and the profound implications of urinary diversion, including lifestyle changes, body image, and sexual function.
  
  \item **Preoperative Antibiotics:** Intravenous broad-spectrum antibiotics are administered shortly before incision to reduce the risk of surgical site infection.
\end{itemize}

\textbf{Factors requiring optimization, alternative treatment, or a change in intent:}
\begin{itemize}
  \item Metastatic disease beyond a potentially curable setting may shift treatment toward systemic or palliative therapy, but selected metastasectomy or consolidative strategies require multidisciplinary assessment.
  
  \item Prohibitive surgical risk due to severe, uncorrectable comorbidities (e.g., severe cardiac failure, uncontrolled sepsis, acute myocardial infarction, severe respiratory failure, extreme frailty with ASA physical status V).
  
  \item Severe coagulopathy that cannot be corrected requires postponement when clinically feasible or an individualized urgent plan.
  
  \item Patient refusal after thorough discussion of risks and benefits.
\end{itemize}

\textbf{Situations requiring optimization, staging, or an alternative treatment plan:}
\begin{itemize}
  \item Severe cardiac or pulmonary disease that significantly increases perioperative risk, but may be optimized with medical management and close monitoring.
  
  \item Poor nutritional status or cachexia, which can impair wound healing and increase infection risk. Preoperative nutritional support and optimization are crucial.
  
  \item Extensive prior abdominal surgery or radiation therapy to the pelvis, which can lead to dense adhesions, increase operative difficulty, and elevate the risk of bowel injury or fistula formation.
  
  \item Inflammatory bowel disease (e.g., Crohn's disease, ulcerative colitis), which may contraindicate the use of affected bowel segments for diversion due to increased risk of complications.
  
  \item Severe renal insufficiency, especially for neobladder formation, due to the increased risk of metabolic acidosis and electrolyte imbalances from urine reabsorption by bowel mucosa.
  
  \item Significant cognitive impairment or limited manual dexterity may affect diversion choice and support needs; it is not an absolute contraindication to cystectomy.
  
  \item Active infection (e.g., urinary tract infection or pneumonia) should be treated and optimized before elective surgery when possible.
  
  \item Morbid obesity, which can increase technical difficulty, operative time, and postoperative complication rates.
\end{itemize}

\section*{Surgical Technique \& Procedure}
The procedure is performed under general anesthesia, often supplemented with epidural analgesia for enhanced postoperative pain control. The patient is typically positioned in a supine or modified dorsal lithotomy position, depending on the surgeon's preference and the need for urethral access for neobladder reconstruction. The surgical approach can be open (via a midline infraumbilical incision extending to the pubic symphysis) or minimally invasive (laparoscopic or robotic-assisted, using several small incisions for ports), with the choice influenced by tumor characteristics, patient anatomy, and surgeon expertise.

Key operative steps include:

\begin{itemize}
  \item **Incision and Access:** A lower midline abdominal incision or multiple robotic ports are placed to gain access to the peritoneal cavity. The bowel is typically packed away from the pelvis to optimize exposure.
  
  \item **Mobilization of the Bladder:** The bladder is carefully dissected from surrounding structures, including the peritoneum, anterior abdominal wall, and pelvic sidewalls. The ureters are identified, mobilized, and transected close to the bladder, with their ends often marked for later identification during diversion.
  
  \item **Pelvic Lymph Node Dissection (PLND):** A meticulous dissection of the regional lymph nodes is performed. The extent of PLND varies but typically includes nodes along the external iliac, internal iliac, and obturator vessels, often extending to the common iliac bifurcation, to ensure adequate oncologic staging and removal of potential micrometastases.
  
  \item **Bladder Removal (En Bloc Resection):**
    \begin{itemize}
      \item In males: The bladder is removed along with the prostate, seminal vesicles, and often a portion of the vas deferens. Nerve-sparing techniques may be attempted to preserve erectile function in select cases.
      \item In females: The bladder is removed with the uterus, fallopian tubes, ovaries, and a portion of the anterior vaginal wall, depending on tumor extent and patient factors.
    \end{itemize}
  
  \item **Urinary Diversion Construction:** A segment of small bowel (typically 15-20 cm of ileum) is isolated.
    \begin{itemize}
      \item For an ileal conduit: The ureters are anastomosed to one end of the isolated ileal segment (Bricker anastomosis), and the other end is brought through the abdominal wall to form a stoma, typically in the right lower quadrant.
      \item For an orthotopic neobladder: A longer segment of bowel (e.g., 40-60 cm of ileum or ileocecal segment) is detubularized and refashioned into a spherical reservoir, which is then anastomosed to the urethra. The ureters are implanted into the neobladder using a reflux-preventing technique.
    \end{itemize}
  
  \item **Bowel Anastomosis:** The continuity of the remaining small bowel is restored with an end-to-end or side-to-side anastomosis, ensuring adequate blood supply and tension-free connection.
  
  \item **Drainage and Closure:** Pelvic drains (e.g., Jackson-Pratt drains) are typically placed to manage fluid collections and monitor for leaks. Ureteral stents may be placed to protect the uretero-enteric anastomoses. The abdominal incision is closed in layers, and the stoma is matured if an incontinent diversion was performed.
\end{itemize}

\section*{Complications \& Risks}
Cystectomy is a major surgical procedure associated with a range of potential risks, which can be broadly categorized as general surgical risks and procedure-specific complications. The overall complication rate can be substantial, reflecting the complexity of the surgery.

\textbf{General Surgical Risks:}
\begin{itemize}
  \item **Bleeding:** Significant intraoperative blood loss requiring transfusion, or postoperative hemorrhage, which may necessitate re-exploration.
  
  \item **Infection:** Surgical site infection (SSI), urinary tract infection (UTI), urosepsis, pneumonia, or other systemic infections (e.g., Clostridioides difficile colitis).
  
  \item **Anesthesia Complications:** Adverse reactions to anesthetic agents, cardiac events (myocardial infarction, arrhythmia), pulmonary complications (atelectasis, pulmonary embolism), stroke, or peripheral nerve injury from positioning.
  
  \item **Deep Vein Thrombosis (DVT) and Pulmonary Embolism (PE):** Blood clot formation in the legs that can travel to the lungs, a potentially life-threatening complication.
  
  \item **Organ Injury:** Accidental injury to adjacent organs such as the bowel, rectum, or major blood vessels during dissection.
\end{itemize}

\textbf{Procedure-Specific Risks:}
\begin{itemize}
  \item **Bowel Complications:**
    \begin{itemize}
      \item Prolonged ileus or bowel obstruction, often requiring nasogastric decompression and delayed oral intake.
      \item Bowel anastomotic leak, leading to peritonitis, intra-abdominal abscess, or fistula formation, requiring reoperation.
      \item Injury to the bowel during dissection, potentially requiring repair or resection.
    \end{itemize}
  
  \item **Urinary Diversion Complications:**
    \begin{itemize}
      \item Uretero-enteric anastomotic leak or stricture, potentially leading to hydronephrosis, renal damage, or urosepsis.
      \item Urinary tract infection, pyelonephritis, or urosepsis, particularly common with diversions due to stasis or reflux.
      \item Metabolic acidosis (especially with ileal conduits or neobladders due to urine reabsorption by bowel mucosa), requiring bicarbonate supplementation.
      \item Stone formation within the diversion or upper urinary tract due to chronic infection or metabolic derangements.
      \item Vitamin B12 deficiency due to resection of the terminal ileum, requiring supplementation.
    \end{itemize}
  
  \item **Stoma-Related Complications (for ileal conduit/continent cutaneous diversions):**
    \begin{itemize}
      \item Stoma stenosis, retraction, necrosis, or prolapse, often requiring revision.
      \item Parastomal hernia, requiring surgical repair.
      \item Skin irritation or breakdown around the stoma due to urine leakage or improper appliance fit.
    \end{itemize}
  
  \item **Neobladder-Specific Complications:**
    \begin{itemize}
      \item Urinary incontinence (stress or nocturnal), requiring pelvic floor therapy or further intervention.
      \item Difficulty voiding, requiring intermittent catheterization due to poor contractility or outflow obstruction.
      \item Neobladder-urethral anastomotic stricture, requiring endoscopic or surgical dilation.
      \item Mucus production, which can cause obstruction or require irrigation.
    \end{itemize}
  
  \item **Sexual Dysfunction:**
    \begin{itemize}
      \item Erectile dysfunction in men due to nerve injury, even with nerve-sparing techniques, often requiring medical or surgical management.
      \item Vaginal shortening or dyspareunia in women, impacting sexual quality of life.
    \end{itemize}
  
  \item **Wound Complications:** Wound dehiscence, seroma, hematoma, or incisional hernia formation at the abdominal incision site.
  
  \item **Lymphocele:** Collection of lymphatic fluid in the pelvis after lymph node dissection, which may require drainage.
  
  \item **Urethral Recurrence:** Although rare, cancer can recur in the remaining urethra, especially in patients with a history of carcinoma in situ.
\end{itemize}

\section*{Postoperative Care \& Recovery}
Immediately after cystectomy, patients are transferred to the Post-Anesthesia Care Unit (PACU) for close monitoring of vital signs, pain control, and fluid balance. Pain management is critical and often involves patient-controlled analgesia (PCA) or epidural catheters. A nasogastric tube (NGT) is typically in place to decompress the stomach and prevent postoperative ileus. Early mobilization, often within 24 hours, is strongly encouraged to prevent complications like deep vein thrombosis, pulmonary embolism, and pneumonia. Intravenous fluids are administered to maintain hydration and electrolyte balance.

The initial hospitalization usually lasts 7 to 14 days, depending on the patient's recovery trajectory and the type of urinary diversion. Within the first 24-48 hours, the NGT may be removed if bowel function shows signs of recovery (e.g., flatus, bowel sounds), and a clear liquid diet is gradually advanced to a regular diet as tolerated. Patients with an ileal conduit will begin learning stoma care and appliance management from an ostomy nurse, while those with a neobladder will start bladder training and potentially self-catheterization if needed. Pelvic drains are monitored for output and typically removed when drainage is minimal (e.g., <50 mL/day). Ureteral stents, if placed, are usually removed several weeks post-discharge. Full recovery from such a major surgery can take several weeks to months, during which patients gradually regain strength, energy, and adapt to their new urinary diversion, requiring significant physical and psychological adjustment.

During radical cystectomy, the surgeon meticulously documents intraoperative findings, including the size, location, and extent of the bladder tumor, its relationship to surrounding structures, and the presence of any suspicious lymph nodes or peritoneal implants. The appearance of the bladder wall, ureteral orifices, and adjacent organs (prostate, seminal vesicles, uterus, vagina) is noted. Any unexpected findings, such as dense adhesions from prior surgery or inflammation, are also recorded. The resected bladder, prostate/seminal vesicles (in males), or uterus/vagina (in females), along with all dissected lymph nodes, are sent to pathology for detailed examination.

The final pathology report is the most critical "result" of a cystectomy. It provides definitive information on the tumor's stage (pTNM classification), histological type (e.g., urothelial carcinoma, squamous cell carcinoma), grade, presence of lymphovascular invasion, and most importantly, the status of surgical margins (whether tumor cells are present at the edges of the resected tissue). It also details the number of lymph nodes removed and how many, if any, contain cancer cells. This comprehensive report guides subsequent oncologic management and surveillance strategies, helping the care team determine the need for further adjuvant therapies (e.g., chemotherapy, radiation) to reduce the risk of cancer recurrence.

A normal, successful postoperative course following cystectomy is characterized by several key outcomes. Clinically, this translates to a smooth recovery without major complications such as significant bleeding, infection, or anastomotic leaks. Effective pain management allows for early mobilization and participation in physical therapy. Bowel function typically returns within 3-7 days, allowing for advancement of diet and removal of the nasogastric tube. Patients should demonstrate progressive improvement in strength and energy levels, with a gradual return to baseline activities.

For patients with an ileal conduit, a normal course involves a well-healed stoma that drains urine effectively into an appliance, with the patient demonstrating proficiency in stoma care and management. For those with a neobladder, normal results include successful urethral voiding, achieving continence (especially daytime), and the absence of significant post-void residual urine, often after a period of bladder training. Ultimately, a normal expected course signifies the resolution of the original symptoms or disease indication, preservation of renal function, and a return to an acceptable quality of life, allowing the patient to resume daily activities and engage in long-term oncologic surveillance.

Post-cystectomy follow-up is multifaceted, focusing on oncologic surveillance, management of the urinary diversion, and overall patient well-being.

\begin{itemize}
  \item **Initial Postoperative Visits:** Typically within 2-4 weeks for wound check, drain removal (if still present), and initial discussion of pathology results and future treatment plans.
  
  \item **Stent Removal:** Ureteral stents, if placed, are usually removed endoscopically 3-6 weeks post-surgery, often in an outpatient setting.
  
  \item **Oncologic Surveillance:** Regular imaging (CT scan of chest, abdomen, and pelvis) every 3-6 months for the first 2-3 years, then annually, to monitor for local recurrence or distant metastases. Cystoscopy of the urethra may be performed if the urethra was spared.
  
  \item **Laboratory Monitoring:** Periodic blood tests to assess renal function (creatinine, electrolytes), vitamin B12 levels (due to ileal resection), and acid-base balance (especially with neobladders, to monitor for metabolic acidosis).
  
  \item **Diversion-Specific Monitoring:**
    \begin{itemize}
      \item For ileal conduits: Annual renal ultrasound to check for hydronephrosis, and occasional stoma evaluation by an ostomy nurse.
      \item For neobladders: Regular assessment of voiding patterns, post-void residuals, and potentially cystoscopy of the neobladder to check for stones or tumor recurrence.
    \end{itemize}
  
  \item **Urine Cytology:** May be performed periodically, especially for neobladder patients, to screen for upper tract recurrence.
  
  \item **Rehabilitation and Support:** Referral to physical therapy for core strengthening, pelvic floor therapy for neobladder patients (to improve continence), and support groups for ostomy care or sexual health counseling.
  
  \item **Activity Resumption:** Gradual return to normal activities, with restrictions on heavy lifting for several weeks to months to prevent incisional hernia.
\end{itemize}

Patients are educated on warning signs such as fever, flank pain, changes in urine output, persistent nausea/vomiting, new masses, or unexplained weight loss, which warrant immediate medical attention.

\section*{Outcomes \& Prognosis}
Bladder cancer is the 4th most common cancer in men and the 9th most common cancer overall in the United States, with an estimated 83,730 new cases and 17,200 deaths annually. It predominantly affects older adults, with the median age at diagnosis being 73 years, and has a significant male predominance (approximately 4 times more common in men than women). Smoking is the most significant risk factor, accounting for about half of all bladder cancer cases, followed by occupational exposure to certain chemicals (e.g., aromatic amines, dyes, rubber). While most bladder cancers are non-muscle invasive at diagnosis, approximately 20-25\% present as muscle-invasive bladder cancer (MIBC) or progress to MIBC, necessitating radical cystectomy. The incidence of MIBC requiring cystectomy is therefore a substantial public health concern, with associated morbidity rates ranging from 30-60\% for any complication and mortality rates typically between 1-3\% in high-volume centers, though it can be higher in less experienced settings or in patients with significant comorbidities.

The outcomes and prognosis following radical cystectomy are primarily dictated by the pathological stage of the bladder cancer, particularly the depth of invasion (pT stage), lymph node status (pN stage), and surgical margin status. For patients with organ-confined muscle-invasive bladder cancer (pT2 N0 M0), 5-year survival rates can range from 60-80\%. However, if lymph nodes are positive (pN+), the 5-year survival rate drops significantly to 20-40\%, highlighting the importance of comprehensive lymph node dissection. Positive surgical margins are a strong predictor of local recurrence and worse overall survival. The choice of urinary diversion generally does not impact oncologic outcomes, but it significantly affects functional outcomes and quality of life. Orthotopic neobladders offer the potential for improved body image and sexual function compared to ileal conduits, though they carry a higher risk of urinary incontinence and the need for intermittent catheterization. Recurrence patterns include local recurrence in the pelvis or urethra, and distant metastases, most commonly to the lungs, liver, and bone. Prognostic factors also include tumor grade, presence of lymphovascular invasion, and patient comorbidities. Adjuvant chemotherapy for high-risk patients (e.g., pT3/T4a or pN+) has been shown in trials to improve survival outcomes, while neoadjuvant chemotherapy is now standard of care for eligible MIBC patients.

\section*{References}
\begin{enumerate}
  \item MedlinePlus. Cystectomy. U.S. National Library of Medicine; 2025. Available from: \url{https://medlineplus.gov/ency/article/002940.htm}
  
  \item Wein AJ, Kavoussi LR, Partin AW, Peters CA, editors. Campbell-Walsh-Wein Urology. 12th ed. Philadelphia, PA: Elsevier; 2025.
  
  \item National Comprehensive Cancer Network (NCCN). NCCN Clinical Practice Guidelines in Oncology: Bladder Cancer. Version 1.2025. Available from: \url{https://www.nccn.org/professionals/physician_gls/pdf/bladder.pdf}
  
  \item Bricker EM. Bladder substitution with a segment of ileum. Ann Surg. 1950;132(3):372-381.
\end{enumerate}

\clearpage
